\documentclass[a4paper,11pt]{article}

\usepackage[utf8]{inputenc}
\usepackage[T1]{fontenc}
\usepackage[french]{babel}
\usepackage{lmodern}
\usepackage{microtype}
\usepackage{amsmath}
\usepackage{booktabs}
\usepackage{longtable}
\usepackage{array}
\usepackage{geometry}
\usepackage{xcolor}
\usepackage{hyperref}

\geometry{margin=2.5cm}
\hypersetup{
  colorlinks=true,
  linkcolor=blue!45!black,
  urlcolor=blue!45!black,
  citecolor=blue!45!black,
  pdftitle={ASTRA P79 Results},
  pdfauthor={ASTRA / Representations Procedurales Adressables}
}

\setlength{\parindent}{0pt}
\setlength{\parskip}{0.60em}
\renewcommand{\arraystretch}{1.12}
\setlength{\emergencystretch}{3em}
\sloppy

\newcommand{\Metric}[2]{#1 & #2\\}
\newcommand{\Decision}{\texttt{RECALIBRATE\_P79\_LEVEL1\_ROUTER\_POLICY}}

\title{\textbf{ASTRA-P79 Results}\\Level-1 Address Space Router}
\author{ASTRA / Repr\'esentations Proc\'edurales Adressables}
\date{2026-05-01}

\begin{document}
\maketitle

\begin{abstract}
P79 teste un routeur d'espace d'adressage de niveau 1. Le routeur choisit
localement entre trie de chemins, DAG de contenu, espace produit type, graphe,
hybride multi-index et baselines grille. Le pipeline reste living-memory:
encode, open, address lookup, read/query, update/delete, audit, compact,
close/reopen et verification du guard. Le routeur economise des bytes d'index
et ameliore le cout d'adressage face a l'hybride seul, mais il rate de peu le
seuil de promotion ratio/hybrid.
\end{abstract}

\section{Objectif}

P78 a montre que l'espace hybride multi-index maximise le ratio living, tandis
que le trie de chemins garde le meilleur cout d'adressage. P79 teste donc un
routeur niveau 1 au lieu d'une topologie unique globale.

L'espace virtuel procedural reste local: il se construit quand une adresse est
atteinte. Les bytes virtuels reportes sont des equivalents de materialisation,
pas des bytes stockes.

\section{Syntaxe Atlas}

P79 ajoute une syntaxe declarative specialisee:

\begin{itemize}
  \item \texttt{p79\_level1\_router\_probe};
  \item \texttt{level1\_router};
  \item \texttt{level1\_router\_gates}.
\end{itemize}

Les gates exigent living-memory, ratio living primaire, metriques d'espace
virtuel, construction locale a l'adresse, guard sans faux gain, index niveau 1
non cache, lookup borne, ratio routeur/oracle minimal et bytes virtuels comme
equivalents.

\section{Routeur et oracle}

La policy \texttt{p79-router} route:
\begin{itemize}
  \item chemins, modules et JSON paths vers le trie;
  \item chunks binaires vers le DAG de contenu;
  \item namespaces types vers l'espace produit;
  \item relations vers le graphe;
  \item acces multiples vers l'hybride;
  \item projections regulieres vers la grille 3D;
  \item guard vers refus ou raw no-gain.
\end{itemize}

L'oracle compare la selection du routeur a la meilleure topologie observee par
groupe de features.

\section{Campagnes}

\begin{longtable}{@{}p{0.55\linewidth}r@{}}
\toprule
Metrique & Valeur\\
\midrule
\Metric{target source bytes}{10,485,760}
\Metric{actual source bytes}{10,485,760}
\Metric{cycles standard / ambitious}{10 / 25}
\Metric{queries standard / ambitious}{10,000 / 50,000}
\Metric{updates standard / ambitious}{1,000 / 5,000}
\Metric{deletes standard / ambitious}{100 / 500}
\Metric{compact}{adaptive}
\Metric{adaptive}{on}
\bottomrule
\end{longtable}

\section{Espace virtuel}

\begin{longtable}{@{}p{0.55\linewidth}r@{}}
\toprule
Metrique & Valeur\\
\midrule
\Metric{effective addresses}{222,464,000,000}
\Metric{virtual cells}{10,000}
\Metric{virtual fibers}{40,000}
\Metric{virtual chunks}{40,000}
\Metric{virtual effective bytes equivalent}{96,415,629}
\Metric{limiting factor}{index size}
\Metric{bytes are equivalent, not stored}{true}
\bottomrule
\end{longtable}

\section{Resultats standard}

\begin{longtable}{@{}p{0.55\linewidth}r@{}}
\toprule
Metrique & Valeur\\
\midrule
\Metric{best single topology}{hybrid multi-index}
\Metric{ratio living router}{5.174901}
\Metric{ratio living hybrid only}{5.340001}
\Metric{ratio living oracle}{5.402000}
\Metric{router / hybrid}{0.969082}
\Metric{router / oracle}{0.957960}
\Metric{lookup p95 router}{7.800}
\Metric{lookup p95 path trie}{7.000}
\Metric{lookup p95 bytes read}{3,840}
\Metric{index bytes router}{1,080,033}
\Metric{index bytes hybrid}{1,342,177}
\Metric{index saved vs hybrid}{262,144}
\Metric{wrong routes}{110}
\Metric{wrong route cost}{359}
\Metric{CRUD success rate}{1.000000}
\Metric{retrieval success rate}{1.000000}
\Metric{reopen equivalence}{true}
\Metric{drift status}{NO\_DRIFT}
\Metric{guard decision}{NO\_GO guard refused}
\bottomrule
\end{longtable}

\section{Resultats ambitious}

La campagne ambitious a ete executee. Les ratios et les couts d'adressage sont
stables; le wrong route cost monte a 407 sous pression d'update plus forte. La
decision ne change pas.

\section{Phase map}

\begin{longtable}{@{}p{0.55\linewidth}r@{}}
\toprule
Statut & Count\\
\midrule
\Metric{green promote}{82}
\Metric{yellow recalibrate}{46}
\Metric{red no-go}{11}
\Metric{grey not tested}{0}
\bottomrule
\end{longtable}

Les rouges viennent surtout des cas ou le routeur choisit une topologie trop
simple pour des features relationnelles ou des projections regulieres.

\section{Decision}

\Decision

Le routeur satisfait retrieval, CRUD, reopen, guard et drift. Il economise
262,144 bytes d'index face a hybrid-only et reste dans l'enveloppe lookup du
path-trie: \texttt{7.800 <= 7.000 * 1.15}. Il rate cependant le seuil
ratio/hybrid: \texttt{0.969082 < 0.970000}.

\section{Limites}

Le routeur est deterministe et compact. Il ne pretend pas stocker plus
d'information par bit. Il choisit seulement une topologie d'adressage locale
pour reduire index et lookup quand la structure de la donnee s'y prete.

\section{Recommandation P80}

P80 devrait calibrer le routeur niveau 1 contre l'oracle: corriger les mauvais
routages relation-heavy, logs sous haute pression et projections CSV, tout en
gardant le gain d'index et en franchissant le seuil \texttt{0.97 * hybrid}.

\end{document}
